%% ED26-P08-inquisitive-0001; Ciardelli, Inquisitive Logic, Chapters 2--4.
\chapter{Inquisitive Propositional Logic}
\label{sec:p08-inquisitive-propositional}

An answer to one question can determine an answer to another without telling
us in advance which answer will result. Inquisitive propositional logic
\(\mathsf{InqB}\) represents this form of inference in the same language as
ordinary propositional statements. Its semantic relation is support by an
information state. Its proofs can transform supplied resolutions of their
premises into a resolution of their conclusion.

This treatment reconstructs Chapters~2--4 of Ciardelli's
\emph{Inquisitive Logic: Consequence and Inference in the Realm of Questions}
\cite{Ciardelli2023InquisitivePropositional}. The presentation and examples
are adapted under the source's Creative Commons Attribution 4.0 licence
\url{https://creativecommons.org/licenses/by/4.0/}; the organization,
wording and expanded derivations here differ from the source.
\index{inquisitive logic}\index{question entailment}\index{information state}

\section{Information, support and dependence}

\begin{hbdefinition}{def:p08-inq-states}{Information states and enhancement}
In a space \(W\) of possible worlds, an \emph{information state} is a subset
\(s\subseteq W\). A smaller state contains at least as much information:
\(t\subseteq s\) is an \emph{enhancement} of \(s\). The empty state is
inconsistent information; a nonempty state is consistent. Information states
here represent sets of possibilities, rather than only the states an agent
can physically or cognitively reach.
\end{hbdefinition}

To support a statement is to establish that it is true. To support a question
is to settle it. A single state can settle a question in more than one way;
different states can settle the same question by establishing different
answers. In particular, supporting a question does not require endorsing
one answer in advance of receiving the relevant information.

\begin{hbdefinition}{def:p08-inq-propositions}{Inquisitive propositions, generators and alternatives}
An \emph{inquisitive proposition} over \(W\) is a nonempty downward-closed
family \(\mathcal A\subseteq\mathcal P(W)\). For a family \(T\) of states put
\[
 T^\downarrow=\{s\subseteq W:\text{some }a\in T\text{ satisfies }s\subseteq a\}.
\]
A \emph{generator} for \(\mathcal A\) is a family \(T\) with
\(\mathcal A=T^\downarrow\). Its \emph{alternatives}
\(\operatorname{Alt}(\mathcal A)\) are its inclusion-maximal members.
It is \emph{normal} if
\(\mathcal A=\operatorname{Alt}(\mathcal A)^\downarrow\).
A generator is \emph{minimal} if no proper subfamily generates the same
proposition, and \emph{independent} if no two distinct members are comparable
by inclusion. A minimal generator is called a \emph{basis}.
\end{hbdefinition}

\begin{hblemma}{thm:p08-inq-bases}{proved=annex:p8.inq.bases}{Generators, alternatives and bases}
A generator is minimal exactly when it is independent. An inquisitive
proposition has a basis exactly when it is normal; in that case its unique
basis is its family of alternatives. Normality is not automatic for arbitrary
inquisitive propositions. For example, all finite subsets of an infinite
\(W\) form a nonempty downward-closed family with no maximal member.
\end{hblemma}

\begin{hbdefinition}{def:p08-inq-general-entailment}{Support consequence and contextual consequence}
Suppose each sentence \(\varphi\) is assigned a persistent support relation,
with \(\varnothing\) supporting every sentence. Its proposition in a model
\(M\) is \([\varphi]_M=\{s:M,s\models\varphi\}\). Define
\[
 \Phi\models\psi
 \quad\Longleftrightarrow\quad
 \text{in every model and state, }s\models\Phi\text{ implies }s\models\psi,
\]
where \(s\models\Phi\) means support for every member. For a fixed state
\(s\) in \(M\), define \(\Phi\models_s\psi\) by quantifying only over
\(t\subseteq s\). Write \(\models_M\) for \(\models_W\).
Equivalence means mutual support entailment; validity means entailment from
no premises. In this chapter \(\equiv_{\mathrm{InqB}}\) denotes equivalence
in the specific logic defined below.
For the general statement cases, the truth-support bridge means that
\(s\models\alpha\) exactly when \(s\subseteq|\alpha|_M\).
\end{hbdefinition}

There are four useful readings. Statement-to-statement consequence preserves
ordinary truth. A statement entails a question when it resolves it. A
question entails a statement when that statement is presupposed by every
resolution. A question entails another question when resolving the former
determines a resolution of the latter. These are support consequences;
causal determination, explanatory dependence and an agent's ability to
perform an experiment are additional notions.

\begin{hblemma}{thm:p08-inq-generator-entailment}{proved=annex:p8.inq.generator-entailment}{Entailment between information types}
In a fixed model choose generators \(T_i\) for \([\varphi_i]_M\) and \(U\)
for \([\psi]_M\). Then, for a finite list of premises,
\[
 \varphi_1,\ldots,\varphi_n\models_M\psi
 \quad\Longleftrightarrow\quad
 \text{for every }a_i\in T_i\text{ there is }b\in U
 \text{ with }\bigcap_{i=1}^n a_i\subseteq b.
\]
The empty intersection is \(W\). For a statement \(\alpha\), the generator
is the singleton \(\{|\alpha|_M\}\), where \(|\alpha|_M\) is its truth set.
Thus a statement resolves a question precisely when its truth set is included
in one generating resolution; a question entails a statement precisely when
every generating resolution is included in the statement's truth set.
\end{hblemma}

\begin{hbexample}{def:p08-inq-die-information}{Parity, range and outcome}
Let \(W=\{1,2,3,4,5,6\}\), with one world for each possible die outcome.
The parity question has generators \(\{1,3,5\}\) and \(\{2,4,6\}\).
The range question has generators \(\{1,2\},\{3,4\},\{5,6\}\).
The outcome question has the six singleton generators. Intersecting a parity
alternative with a range alternative gives one singleton, so parity and range
jointly determine outcome. Range alone does not: \(\{1,2\}\) resolves range
but leaves outcome open. Under the information \(\{2,3,5\}\) that the outcome
is prime, each range alternative contains at most one live world, so range
then determines outcome.
\end{hbexample}

\section{The propositional language and its support semantics}

\begin{hbdefinition}{def:p08-inq-language}{Classical formulas and \(\mathsf{InqB}\) formulas}
Fix a countably infinite set \(P\) of atoms. The classical fragment is
\[
 \alpha ::=p\mid\bot\mid(\alpha\land\alpha)\mid(\alpha\to\alpha).
\]
The full language additionally permits inquisitive disjunction, printed here
as the doubled-vee connective \(\mathbin{\vee\!\vee}\):
\[
 \varphi ::=p\mid\bot\mid(\varphi\land\varphi)\mid
 (\varphi\to\varphi)\mid(\varphi\mathbin{\vee\!\vee}\varphi).
\]
The defined connectives are
\[
\begin{aligned}
 \neg\varphi&:=\varphi\to\bot,\qquad \top:=\neg\bot,\\
 \varphi\lor\psi&:=\neg(\neg\varphi\land\neg\psi),\\
 \varphi\leftrightarrow\psi&:=(\varphi\to\psi)\land(\psi\to\varphi),\\
 ?\varphi&:=\varphi\mathbin{\vee\!\vee}\neg\varphi.
\end{aligned}
\]
In rule side conditions, \(\alpha,\beta\) mean syntactically classical
formulas; \(\varphi,\psi,\chi\) may be arbitrary formulas. A finite iterated
inquisitive disjunction is a bracketing of the binary connective; its empty
instance denotes \(\bot\), and an empty conjunction denotes \(\top\).
\end{hbdefinition}

\begin{hbdefinition}{def:p08-inq-support}{Models and all support clauses}
A propositional information model is \(M=(W,V)\), where \(W\) is a set and
\(V:W\times P\to\{0,1\}\). The set \(W\) may be empty. For \(s\subseteq W\),
\[
\begin{aligned}
 M,s\models p&\quad\Longleftrightarrow\quad V(w,p)=1\text{ for every }w\in s,\\
 M,s\models\bot&\quad\Longleftrightarrow\quad s=\varnothing,\\
 M,s\models\varphi\land\psi&\quad\Longleftrightarrow\quad
 M,s\models\varphi\text{ and }M,s\models\psi,\\
 M,s\models\varphi\mathbin{\vee\!\vee}\psi&\quad\Longleftrightarrow\quad
 M,s\models\varphi\text{ or }M,s\models\psi,\\
 M,s\models\varphi\to\psi&\quad\Longleftrightarrow\quad
 \text{for every }t\subseteq s,
 \ M,t\models\varphi\text{ implies }M,t\models\psi.
\end{aligned}
\]
Truth at a world is the derived notion \(M,w\models\varphi\) meaning
\(M,\{w\}\models\varphi\). Put
\(|\varphi|_M=\{w:M,w\models\varphi\}\).
To \emph{rule out} \(\varphi\) at \(s\) means that no nonempty
\(t\subseteq s\) supports it.
\end{hbdefinition}

\begin{hblemma}{thm:p08-inq-basic}{proved=annex:p8.inq.basic}{Persistence, the empty state, locality and truth}
Every formula is supported at \(\varnothing\), and support persists to subsets.
Restricting a model to \(s\) does not change support at any \(t\subseteq s\).
At singleton states all connectives have their ordinary classical truth
conditions, with both disjunctions giving inclusive disjunction.
Moreover,
\[
 |\varphi|_M=\bigcup[\varphi]_M,
 \qquad
 s\models\neg\varphi\ \Longleftrightarrow\ s\cap|\varphi|_M=\varnothing.
\]
Thus negation expresses ruling out, rather than failure of support. In
particular, \(s\not\models\varphi\) does not imply
\(s\models\neg\varphi\).
\end{hblemma}

\begin{hbdefinition}{def:p08-inq-flatness}{Statements, questions and the classical variant}
A formula is \emph{truth-conditional}, or \emph{flat}, if in every model and state
\(s\models\varphi\) exactly when every \(w\in s\) satisfies it.
In this presentation a \emph{statement} is a flat formula and a
\emph{question} is a non-flat formula. This semantic terminology differs
from the convention that calls every non-informative formula a question.
The classical variant \(\varphi^{\mathrm{cl}}\) replaces every inquisitive
disjunction by the defined classical disjunction. Its truth set is the
\emph{presupposition} of \(\varphi\).
\end{hbdefinition}

\begin{hblemma}{thm:p08-inq-flatness}{proved=annex:p8.inq.flatness}{Truth-conditionality and double negation}
Classical formulas are flat. A formula is flat exactly when it is equivalent
to its classical variant, equivalently to some classical formula, equivalently
to its double negation. For every \(\varphi\),
\[
 |\varphi|_M=|\varphi^{\mathrm{cl}}|_M,
 \quad\neg\neg\varphi\equiv_{\mathrm{InqB}}\varphi^{\mathrm{cl}},
 \quad\varphi\models\neg\neg\varphi.
\]
Negations and classical disjunctions are always flat, as is an implication
with flat consequent. A conjunction of flat formulas is flat. Flatness is
equivalent to closure of the supporting states under arbitrary unions.
\end{hblemma}

\begin{hbexample}{def:p08-inq-polar-counterstate}{Truth does not resolve a polar question}
Let \(W=\{w_0,w_1\}\), with \(p\) false at \(w_0\) and true at \(w_1\).
The state \(W\) supports \(p\lor\neg p\), but supports neither \(p\) nor
\(\neg p\), and therefore does not support \(?p\). Each singleton supports
\(?p\). Thus singleton truth cannot test inquisitive validity, and
\(\neg\neg?p\models?p\) fails. The negation \(\neg?p\), by contrast,
is equivalent to \(\bot\): ruling out the possibility of resolving a
tautologically presupposed polar question requires the inconsistent state.
\end{hbexample}

\section{Implication, context and presupposition}

\begin{hblemma}{thm:p08-inq-implication}{proved=annex:p8.inq.implication}{Contextual consequence and generalized supposition}
For every state \(s\),
\[
 s\models\varphi\to\psi\quad\Longleftrightarrow\quad\varphi\models_s\psi.
\]
Nested implication internalizes any finite list of premises. If \(\alpha\)
is flat, then
\[
 s\models\alpha\to\psi
 \quad\Longleftrightarrow\quad s\cap|\alpha|_M\models\psi.
\]
More generally, for any generator \(T\) of \([\varphi]_M\),
\[
 s\models\varphi\to\psi
 \quad\Longleftrightarrow\quad
 s\cap a\models\psi\text{ for every }a\in T.
\]
For flat \(\alpha\), replacing an antecedent \(\varphi\) by
\(\varphi^{\mathrm{cl}}\) does not change \(\varphi\to\alpha\).
For a set of premises with flat conclusion, support entailment reduces to
worldwise truth entailment, and every premise may be replaced by its classical
variant. In particular, a question entails precisely the statements entailed
by its presupposition.
\end{hblemma}

\begin{hbexample}{def:p08-inq-dependence-functions}{Four ways for whether-\(p\) to determine whether-\(q\)}
In a state \(s\), \(?p\to?q\) is supported exactly when the truth value of
\(q\) is a function of the truth value of \(p\) across \(s\). The four
possible Boolean functions give the classical conditions
\(q,\neg q,p\leftrightarrow q,p\leftrightarrow\neg q\).
Their truth sets generate the implication's support. In the full four-world
model these are four different two-world alternatives, including both
diagonals. The implication expresses dependence relative to a state; a
statement that an agent knows this dependence requires an additional modal
language.
\end{hbexample}

\section{Resolutions and normal form}

\begin{hbdefinition}{def:p08-inq-resolutions}{Recursive resolutions}
Fix a syntactic ordering of classical formulas, used to order finite
conjunctions. Define the finite nonempty set \(R(\varphi)\) by
\[
\begin{aligned}
 R(p)&=\{p\},&R(\bot)&=\{\bot\},\\
 R(\varphi\land\psi)&=\{\alpha\land\beta:
              \alpha\in R(\varphi),\ \beta\in R(\psi)\},\\
 R(\varphi\mathbin{\vee\!\vee}\psi)&=R(\varphi)\cup R(\psi),\\
 R(\varphi\to\psi)&=
 \left\{\gamma_f:f:R(\varphi)\to R(\psi)\right\},
\end{aligned}
\]
where \(\gamma_f=\bigwedge_{\alpha\in R(\varphi)}(\alpha\to f(\alpha))\).
Singleton conjunctions retain their sole formula. A resolution is a
classical formula, not a proof or an empirical observation. For a set
\(\Phi\), let \(R(\Phi)\) be the family of images of selections
\(r:\Phi\to L_c\) with \(r(\varphi)\in R(\varphi)\), where \(L_c\)
is the classical fragment. For a finite list of assumptions, a resolution
is a list choosing one resolution of each listed assumption.
\end{hbdefinition}

\begin{hbtheorem}{thm:p08-inq-resolution-semantics}{proved=annex:p8.inq.resolution-semantics}{Resolution theorem for support and normal form}
Every \(R(\varphi)\) is finite, nonempty and classical, and
\(R(\alpha)=\{\alpha\}\) for classical \(\alpha\). For every model and state,
\[
 s\models\varphi
 \quad\Longleftrightarrow\quad
 s\models\alpha\text{ for some }\alpha\in R(\varphi).
\]
Consequently \(\varphi\) is equivalent to the inquisitive disjunction of
its resolutions, every resolution entails \(\varphi\), and
\([\varphi]_M=\{|\alpha|_M:\alpha\in R(\varphi)\}^\downarrow\).
Every propositional formula is normal. It is flat exactly when it has one
alternative in every model, necessarily its truth set. A state supports a
set \(\Phi\) exactly when it supports some \(\Gamma\in R(\Phi)\).
\end{hbtheorem}

\begin{hbdefinition}{def:p08-inq-logical-dependence}{Dependence functions}
A dependence function from \(\Phi\) to \(\psi\) relative to \(s\) maps
each \(\Gamma\in R(\Phi)\) to \(f(\Gamma)\in R(\psi)\) such that
\(\Gamma\models_s f(\Gamma)\). It is logical if this entailment holds
in every model and state. For a single antecedent use a function
\(f:R(\varphi)\to R(\psi)\); its associated \(\gamma_f\) expresses
that this particular resolution-transfer function works.
\end{hbdefinition}

\begin{hbexample}{def:p08-inq-contextual-function}{One resolution-transfer function in two states}
Use the eight valuations of atoms \(p,q,r\), encoded by three bits in that
order. Define \(f(p)=q\) and \(f(\neg p)=r\). In
\(s_1=\{111,110,011,001\}\), every \(p\)-world makes \(q\) true and
every \(\neg p\)-world makes \(r\) true. Hence \(f\) transfers resolutions
of \(?p\) to resolutions of \(q\mathbin{\vee\!\vee}r\) in \(s_1\).
In \(s_2=\{100,110,011,001\}\), it fails: world 100 satisfies \(p\)
but not \(q\). Thus a proposed resolution-transfer map
requires the stated contextual premises; merely specifying its input
and output sets does not make it a dependence function.
\end{hbexample}

\begin{hbtheorem}{thm:p08-inq-split}{proved=annex:p8.inq.split}{Specificity, split and resolution of entailment}
Let \(\Gamma\) contain flat formulas. Put
\(|\Gamma|_M=\bigcap_{\alpha\in\Gamma}|\alpha|_M\). Then
\[
 \Gamma\models_s\varphi
 \Longleftrightarrow s\cap|\Gamma|_M\models\varphi,
 \qquad
 \Gamma\models\varphi
 \Longleftrightarrow |\Gamma|_M\models\varphi\text{ for every }M.
\]
For either logical or contextual entailment, \(\Gamma\) entails an
inquisitive disjunction exactly when it entails one particular disjunct,
and entails \(\varphi\) exactly when it entails one resolution of it.
For arbitrary \(\Phi\),
\[
 \Phi\models\psi
 \Longleftrightarrow
 \text{for every }\Gamma\in R(\Phi)\text{ some }\beta\in R(\psi)
 \text{ satisfies }\Gamma\models\beta.
\]
The same assertion holds relative to \(s\). Equivalently, the appropriate
dependence function exists. For a single antecedent,
\(s\models\gamma_f\) exactly when the particular map \(f\) works
in \(s\), and \(s\models\varphi\to\psi\) exactly when some
such map works. In particular,
\[
 \alpha\to(\varphi\mathbin{\vee\!\vee}\psi)
 \equiv_{\mathrm{InqB}}
 (\alpha\to\varphi)\mathbin{\vee\!\vee}(\alpha\to\psi)
 \quad(\alpha\text{ flat}).
\]
\end{hbtheorem}

\section{Natural deduction and resolutions of proofs}

\begin{hbdefinition}{def:p08-inq-proofs}{Proofs and assumption discharge}
A proof is a finite natural-deduction tree using the following sheet.
Its undischarged assumptions must be included in the stated premise set.
An implication introduction or disjunction elimination may discharge any
chosen number of occurrences of its displayed temporary assumption,
including zero. The source's classical-formula restrictions apply to the
entire formula in the side condition, not merely to its principal connective.
Write \(P:\Phi\vdash_{\mathrm{InqB}}\psi\) for a particular proof and
\(\Phi\vdash_{\mathrm{InqB}}\psi\) for the existence of one.
\end{hbdefinition}

\showcalculus{p08-inquisitive}

\begin{hbtheorem}{thm:p08-inq-soundness}{proved=annex:p8.inq.soundness}{Soundness and the derived classical rules}
If \(\Phi\vdash_{\mathrm{InqB}}\psi\), then \(\Phi\models\psi\).
The usual negation rules are derived from implication and falsum.
Classical disjunction introduction is derivable for arbitrary operands;
classical disjunction elimination is derivable when its conclusion is
classical. Removing that restriction is unsound. Double-negation elimination
for arbitrary questions is likewise unsound.
\end{hbtheorem}

\begin{hbtheorem}{thm:p08-inq-proof-resolution}{proved=annex:p8.inq.proof-resolution}{The resolution algorithm for a proof}
Given a finite proof \(P:\varphi_1,\ldots,\varphi_n\vdash_{\mathrm{InqB}}\psi\),
one can construct recursively on its proof tree an algorithm \(F_P\).
For every \(\alpha_i\in R(\varphi_i)\), it returns a classical proof
\[
 F_P(\alpha_1,\ldots,\alpha_n):
 \alpha_1,\ldots,\alpha_n\vdash\beta
 \quad\text{with }\beta\in R(\psi).
\]
The algorithm covers each introduction and elimination rule, split,
falsum elimination and classical double-negation elimination. Thus the
proof supplies a logical dependence function, not just an assertion that
one exists. The output resolution may depend on the supplied resolutions.
\end{hbtheorem}

\section{Canonical support model and completeness}

\begin{hblemma}{thm:p08-inq-syntactic-resolutions}{proved=annex:p8.inq.syntactic-resolutions}{Provable normal form and traceable deduction failure}
Every formula is interderivable with its inquisitive resolution normal form.
If \(\Phi,\varphi\not\vdash\psi\), some \(\alpha\in R(\varphi)\)
satisfies \(\Phi,\alpha\not\vdash\psi\). If \(\Phi\not\vdash\psi\),
some \(\Gamma\in R(\Phi)\) still does not derive \(\psi\). Consequently,
\[
 \Phi\vdash\psi
 \Longleftrightarrow
 \text{every }\Gamma\in R(\Phi)\text{ derives some }\beta\in R(\psi).
\]
A set of classical premises derives an inquisitive disjunction only if it
derives a particular disjunct.
\end{hblemma}

\begin{hbdefinition}{def:p08-inq-canonical-model}{Classical theories and the canonical model}
A classical theory is a subset \(\Gamma\subseteq L_c\) closed under
derivation of classical formulas in the displayed calculus. It is complete
if \(\bot\notin\Gamma\) and for each classical \(\alpha\), either
\(\alpha\in\Gamma\) or \(\neg\alpha\in\Gamma\).
The canonical model \(M^c=(W^c,V^c)\) has all complete classical theories
as worlds and \(V^c(\Gamma,p)=1\) exactly when \(p\in\Gamma\).
For a state \(S\subseteq W^c\), let \(\bigcap S\) be the common classical
theory, with \(\bigcap\varnothing=L_c\).
\end{hbdefinition}

\begin{hblemma}{thm:p08-inq-support-lemma}{proved=annex:p8.inq.support-lemma}{Classical completion and the canonical support lemma}
Every consistent set of classical formulas extends to a complete classical
theory, and intersections of classical theories are classical theories.
For every canonical state \(S\) and formula \(\varphi\),
\[
 M^c,S\models\varphi
 \quad\Longleftrightarrow\quad
 \bigcap S\vdash_{\mathrm{InqB}}\varphi.
\]
The singleton instance is the truth lemma, but nonsingleton states are
essential to the general statement.
\end{hblemma}

\begin{hbtheorem}{thm:p08-inq-completeness}{proved=annex:p8.inq.completeness}{Completeness of inquisitive natural deduction}
For any premise set \(\Phi\) and conclusion \(\psi\) in the fixed
countable propositional language,
\[
 \Phi\models\psi\quad\Longleftrightarrow\quad
 \Phi\vdash_{\mathrm{InqB}}\psi.
\]
If the conclusion is not derivable, the canonical model has a finite state
supporting every premise and failing to support the conclusion. It can be
chosen with at most \(|R(\psi)|\) worlds. As a consequence, support entailment
is compact: an entailed conclusion already follows from finitely many premises.
\end{hbtheorem}

\begin{hbexample}{def:p08-inq-opposite-answers}{Opposite answers and a proof of dependence}
Suppose \(p\leftrightarrow\neg q\). Temporarily suppose \(?p\).
In its \(p\) branch, implication elimination gives \(\neg q\), hence \(?q\)
by inquisitive disjunction introduction. In its \(\neg p\) branch,
temporarily assuming \(\neg q\) and using \(\neg q\to p\) gives a
contradiction, so \(\neg\neg q\), and classical double-negation elimination
gives \(q\); again introduce \(?q\). Inquisitive disjunction elimination
discharges the two answer assumptions, and implication introduction
discharges \(?p\). The result is
\[
 p\leftrightarrow\neg q\vdash_{\mathrm{InqB}}?p\to?q.
\]
The proof's resolution algorithm sends the answer \(p\) to \(\neg q\)
and \(\neg p\) to \(q\), retaining the background statement.
\end{hbexample}

\section{Expressiveness and effective finite descriptions}

\begin{hbdefinition}{def:p08-inq-profiles}{Finite profiles and state properties}
For finite \(Q\subseteq P\), put \(\Omega_Q=\{0,1\}^Q\). The profile
of a world is its valuation restricted to \(Q\); the profile of a state
is the set of profiles realized by its worlds. A property of pairs
\((M,s)\) is \(Q\)-determined if equal \(Q\)-profiles always have the
same property, even in different models. It is inquisitive if it is
downward closed in each model and includes every empty state.
For \(v\in\Omega_Q\), let \(\chi_v\) be the conjunction of \(q\) when
\(v(q)=1\) and \(\neg q\) when \(v(q)=0\). For
\(A\subseteq\Omega_Q\), let \(\chi_A=\bigvee_{v\in A}\chi_v\), using
classical disjunction and \(\chi_\varnothing=\bot\).
\end{hbdefinition}

\begin{hbtheorem}{thm:p08-inq-expressive}{proved=annex:p8.inq.expressive}{Expressive completeness, local tabularity and decision}
A state property is definable in \(\mathsf{InqB}\) exactly when it is
inquisitive and determined by finitely many atoms. A nonempty downward-closed
\(\mathcal A\subseteq\mathcal P(\Omega_Q)\) is defined by the finite
inquisitive disjunction of the \(\chi_A\) for \(A\in\mathcal A\).
Consequently there are finitely many formulas up to equivalence over a finite
alphabet. If \(D(m)\) denotes the number of antichains of the subset lattice
of an \(m\)-element set, the number over \(n\) atoms is \(D(2^n)-1\);
exactly \(2^{2^n}\) of these classes are flat.
Over one atom there are exactly five classes, represented by
\(\bot,p,\neg p,?p,\top\); over two atoms there are 167.
The one-atom entailment order is
\(\bot<p,\neg p<?p<\top\), with \(p\) and \(\neg p\) incomparable.

Validity is decidable by computing resolutions and testing whether one
resolution is a classical tautology. Equivalently, evaluate the formula
at the entire standard valuation space over its finitely many occurring
atoms. Finite-premise entailment reduces to validity of an implication
from their conjunction. Neither procedure decides an arbitrarily presented
infinite premise set.
\end{hbtheorem}

\begin{hbtheorem}{thm:p08-inq-substitution}{proved=annex:p8.inq.substitution}{Legitimate substitution and the role of atoms}
Uniform substitution of arbitrary formulas does not preserve
\(\mathsf{InqB}\)-consequence. It does preserve consequence when every
atom is replaced by a flat formula. In particular, classical substitutions
are legitimate. Interpreting atoms instead by arbitrary inquisitive
propositions produces a different logic; the present classical atomic
semantics and its normal-form and proof theorems are not thereby preserved.
\end{hbtheorem}

The distinction reflects the role assigned to an atom: it stands for a
statement, not for an arbitrary question. The counterexample substitutes
\(?p\) for \(p\) in classical double-negation elimination. It is the
same two-world failure already exhibited, not a defect in that restricted
classical rule.
The arbitrary-proposition atomic variant is related in the selected source
to Skvortsov's logic of infinite problems and to Medvedev's finite-problem
tradition. The source attributes a nonfinite-axiomatizability result for
that other system to Shehtman and Skvortsov
\cite[Section~3.8]{Ciardelli2023InquisitivePropositional}.
Its atomic semantics and metatheory are separate from the finite
\(\mathsf{InqB}\) calculus given here.

\section{Eliminating connectives and defining operations}

\begin{hbdefinition}{def:p08-inq-connective-definability}{Elimination and uniform definition}
A connective is \emph{eliminable} if every formula can be replaced by an
equivalent formula not containing it. It is \emph{definable} from the other
connectives if one fixed context in them gives an equivalent replacement
for every tuple of operand formulas. A context may contain additional fixed
atoms besides its designated argument places. A set of connectives is
independent if none is so definable from the others.
\end{hbdefinition}

\begin{hbtheorem}{thm:p08-inq-connectives}{proved=annex:p8.inq.connectives}{Independent operations and eliminable connectives}
In \(\mathsf{InqB}\), conjunction is eliminable, while
\(\bot,\to,\mathbin{\vee\!\vee}\) are not. Nevertheless all four primitive
connectives are independent: conjunction has no uniform definition from
the other three. Thus \(\{\bot,\to,\mathbin{\vee\!\vee}\}\) is the unique
minimal expressively complete subset of the given primitives.

In the notational variant \(\mathsf{InqB}^{\neg}\), take
\(\neg,\land,\to,\mathbin{\vee\!\vee}\) as primitives and define
\(\bot\) using a contradiction. Conjunction and implication are jointly
eliminable, and \(\{\neg,\mathbin{\vee\!\vee}\}\) is the unique minimal
complete subset. All four operations are still independent: in particular,
implication has no uniform definition from
\(\neg,\land,\mathbin{\vee\!\vee}\).
\end{hbtheorem}

Elimination can inspect a whole formula and its resolutions; a uniform
definition must work unchanged when questions replace its argument places.
This explains why an equivalent conjunction-free formula does not provide
a general definition of intersection on inquisitive contents. In a logic
closed under arbitrary uniform substitution these notions coincide; that
substitution premise fails here.

\section{Coherence, partitions and neighboring systems}

\begin{hbdefinition}{def:p08-inq-coherence}{Coherence degree}
For \(n\geq1\), a formula is \(n\)-coherent if support at any state is
equivalent to support at all its substates of size at most \(n\).
Its coherence degree \(d(\varphi)\) is the least such \(n\), when it exists.
For properties of states use the same condition, with empty-state membership
understood. A formula is partitional if its alternatives in every model
are pairwise disjoint and cover the model's universe; when the universe
is empty, its sole supporting state is \(\varnothing\).
\end{hbdefinition}

\begin{hbtheorem}{thm:p08-inq-coherence}{proved=annex:p8.inq.coherence}{Finite coherence degrees and the polar fragment}
Every \(\mathsf{InqB}\) formula has finite coherence degree. Degrees obey
\[
\begin{aligned}
 d(p)=d(\bot)&=1,&
 d(\varphi\land\psi)&\leq\max\{d(\varphi),d(\psi)\},\\
 d(\varphi\to\psi)&\leq d(\psi),&
 d(\varphi\mathbin{\vee\!\vee}\psi)&\leq d(\varphi)+d(\psi).
\end{aligned}
\]
Every positive integer occurs as a degree. The fragment
\(\mathsf{InqC}\), whose primitives are \(\bot,\land,\to,?\), is precisely
expressively complete for the finitely determined 2-coherent state properties.
It is strictly less expressive than \(\mathsf{InqB}\).
Partitional formulas are, up to equivalence, finite conjunctions of polar
questions \(?\alpha_i\) with classical \(\alpha_i\).
Equivalently, their support in each model is agreement of every pair of
worlds in the state under one equivalence relation on the universe.
Any state property characterized by its nonempty substates of size at
most \(n\) is downward closed and includes the empty state.
\end{hbtheorem}

The coherence hierarchy also poses a syntactic presentation problem:
to describe, for a fixed degree bound, a fragment whose formulas express
exactly that class of finitely determined properties. This asks for
restrictions on formulas, beyond the semantic test supplied by a finite
profile table. The explicit 2-coherent polar fragment above is one such
presentation.

\begin{hbdefinition}{def:p08-inq-tensor}{Tensor disjunction as a distinct connective}
An extension \(\mathsf{InqB}^{\otimes}\) adds
\[
 s\models\varphi\otimes\psi
 \quad\Longleftrightarrow\quad
 s=s_1\cup s_2\text{ for some }s_1\models\varphi, s_2\models\psi.
\]
The parts may overlap or be empty. This splits the information state;
inquisitive disjunction instead selects a disjunct supported by the whole
state, and classical disjunction preserves only the disjunction's truth set.
\end{hbdefinition}

\begin{hbtheorem}{thm:p08-inq-tensor}{proved=annex:p8.inq.tensor}{Tensor: elimination without definition}
Tensor preserves persistence, the empty-state property and finite-profile
dependence. Every formula of the extension is equivalent to an
\(\mathsf{InqB}\) formula. Nevertheless tensor has no uniform definition
from the four \(\mathsf{InqB}\) primitives.
\end{hbtheorem}

\begin{hbexample}{def:p08-inq-pair-boundary}{Overlapping answers defeat a pair-only test}
Let worlds be the nine two-digit codes \(ij\), with \(i,j\in\{1,2,3\}\).
The question asking for one digit occurring in the code has three alternatives:
all codes containing 1, all codes containing 2, and all codes containing 3.
They overlap. In the state \(\{12,13,23\}\), no digit is guaranteed to
occur, so the question is unresolved. Every pair in that state has a shared
digit and resolves it. Thus pairwise resolution does not in general imply
resolution of the entire state.
\end{hbexample}

For questions whose nonempty alternatives form a partition of a presupposed
subspace, partition refinement and support entailment agree: take the
downward closure of the blocks, including the empty state. If the partition
family is empty, this gives \(\{\varnothing\}\), not an empty proposition.
The block description can be recovered by maximality. General inquisitive
contents permit overlapping alternatives, so they also cover mention-some
and conditional questions beyond this partition case. The pair counterexample
shows why replacing arbitrary states by pairs loses still more information.

\begin{hbtheorem}{thm:p08-inq-partition-limit}{proved=annex:p8.inq.partition-limit}{A partition cannot supply the general implication}
In the six-world die model, no partition content \(Q\) can satisfy
\[
 \lambda,\mathsf{parity}\models_M\mathsf{outcome}
 \quad\Longleftrightarrow\quad \lambda\models_M Q
\]
for every unique-answer question \(\lambda\). The inquisitive implication
does satisfy this equivalence, but its nine maximal supports pair one even
with one odd world and overlap.
\end{hbtheorem}

\section{Finite task tables and worked source problems}

\begin{hbtheorem}{thm:p08-inq-polar-tasks}{proved=annex:p8.inq.polar-tasks}{Polar questions in a fixed model}
For flat \(\alpha,\beta\) in a model \(M\): \(?\alpha\) is supported
everywhere exactly when \(\alpha\) or \(\neg\alpha\) is supported everywhere.
The polar questions \(?\alpha,?\beta\) are equivalent in \(M\) exactly
when their truth sets agree or are complements. If one polar question
entails another in \(M\), either they are equivalent or the conclusion
is supported everywhere. On a nine-world space there are 256 different
polar contents, including the trivial one.
\end{hbtheorem}

\begin{hbexample}{def:p08-inq-code-table}{Information types for a two-digit code}
For the nine-world code space, write \(R_i=\{i1,i2,i3\}\),
\(C_i=\{1i,2i,3i\}\), and \(A_i=R_i\cup C_i\).
Each entry below lists the maximal supporting states; its subsets also
support the corresponding statement or question.
\begin{longtable}{@{}p{0.44\linewidth}p{0.48\linewidth}@{}}
\caption{Support requirements in the code example}\label{tab:p08-inq-code}\\
\hline Content & Alternatives \\\hline\endfirsthead
\caption[]{Support requirements in the code example (continued)}\\
\hline Content & Alternatives \\\hline\endhead
\hline\multicolumn{2}{r}{\emph{continued on next page}}\\
\endfoot
\endlastfoot
The code is 12 & \(\{12\}\). \\
The first digit is 1 & \(R_1\). \\
The code contains 1 & \(A_1\). \\
If its first digit is 1, its second is 2 & \((W\setminus R_1)\cup\{12\}\). \\
Whether the code is 12 & \(\{12\}\), \(W\setminus\{12\}\). \\
Whether the first digit is 1 & \(R_1\), \(W\setminus R_1\). \\
Whether the first digit is 1 or is 2 & \(R_1,R_2\); the third row is outside the presupposition. \\
Whether the code contains 1 & \(A_1\), \(W\setminus A_1\). \\
What the code is & The nine singletons. \\
What its first digit is & \(R_1,R_2,R_3\). \\
One digit absent from the code & \(W\setminus A_1,W\setminus A_2,W\setminus A_3\). \\
If its first digit is 1, what its second is & \((W\setminus R_1)\cup\{1j\}\), for \(j=1,2,3\). \\
\hline
\end{longtable}
The alternative family \(\{11,22,33\}\) versus its complement asks whether
the digits are equal. The three families with first digit less than, equal
to, or greater than the second ask for their ordering. The five alternatives
\(\{11\},\{12,21\},\{13,22,31\},\{23,32\},\{33\}\) ask for the
sum of the digits. These are different information types on the same worlds.
\end{hbexample}

\begin{hbexample}{def:p08-inq-die-implications}{Calculating conditional information types}
In the die space let \(E=\{2,4,6\}\), \(O=\{1,3,5\}\), and
\(L_1=\{1,2\},L_2=\{3,4\},L_3=\{5,6\}\).
The implication from even to low is a statement with truth set
\(O\cup L_1\). The implication from even to range has alternatives
\(O\cup L_i\), while even-to-parity is valid. Range-to-even is the
statement with truth set \(E\). Parity-to-range has the nine alternatives
consisting of one even and one odd outcome. Range-to-parity has the eight
alternatives choosing exactly one outcome from each \(L_i\); these also
are the alternatives for range-to-outcome. Each calculation follows by
intersecting the state with every antecedent alternative and requiring a
consequent resolution.
\end{hbexample}

\begin{hbexample}{def:p08-inq-die-proof}{A proof using an arbitrary answer}
Introduce atoms \(a_1,\ldots,a_6\), an exhaustive classical disjunction
and pairwise exclusivity assumptions. Put
\(\mathsf{prime}=a_2\lor a_3\lor a_5\),
\(\mathsf{range}=(a_1\lor a_2)\mathbin{\vee\!\vee}
 (a_3\lor a_4)\mathbin{\vee\!\vee}(a_5\lor a_6)\), and let
\(\mathsf{outcome}\) be the inquisitive disjunction of all six atoms.
Under \(\mathsf{prime}\), the three range branches yield respectively
\(a_2,a_3,a_5\) by classical reasoning. Introduce \(\mathsf{outcome}\)
in each branch and eliminate the range disjunction. The proof takes any
answer to range and returns the corresponding answer to outcome; it
does not assume which range answer will be supplied.
\end{hbexample}

\begin{hbtheorem}{thm:p08-inq-exercise-inferences}{proved=annex:p8.inq.exercise-inferences}{Further mixed statement-and-question derivations}
The displayed calculus proves
\[
 ?p\land?q\ \vdash\ ?(p\land q),
\]
\[
 p\to(?q\to?r),\neg q\ \vdash\
 (p\to r)\mathbin{\vee\!\vee}(p\to\neg r),
\]
and
\[
 a\lor b,\ a\to\neg m,\ b\to m\ \vdash\
 ?m\to(a\mathbin{\vee\!\vee}b).
\]
If \(m\) means morning, \(h\) means inside the house, and \(b\) means
the butler is the culprit, the two assumptions
\(m\to(h\to b)\) and \(b\to(\neg h\to\neg m)\) yield
\[
 ?h\ \vdash\ m\to?b
\]
in their presence. The two question answers lead to different classical
conclusions; classical background disjunction is not itself an answered question.
\end{hbtheorem}

\begin{hbexample}{def:p08-inq-entailment-table}{A finite entailment lookup}
In the code model use the local abbreviations: \(a\), the code is 12;
\(b\), its first digit is 1; \(c\), it contains 1; \(f\), whether its
first digit is 1; \(h\), whether it contains 1; \(i\), what the code is;
and \(j\), what its first digit is. A 1 at row \(u\), column \(v\)
means \(u\models_M v\); a 0 means failure.
\begin{center}
\begin{tabular}{c|rrrrrrr}
 &\(a\)&\(b\)&\(c\)&\(f\)&\(h\)&\(i\)&\(j\)\\\hline
\(a\)&1&1&1&1&1&1&1\\
\(b\)&0&1&1&1&1&0&1\\
\(c\)&0&0&1&0&1&0&0\\
\(f\)&0&0&0&1&0&0&0\\
\(h\)&0&0&0&0&1&0&0\\
\(i\)&0&0&0&1&1&1&1\\
\(j\)&0&0&0&1&0&0&1
\end{tabular}
\end{center}
Every entry is obtained by checking whether each row alternative is
included in a column alternative. Thus this table does not mistake the
truth of a presupposition for resolution of its question.
\end{hbexample}

\begin{hbexample}{def:p08-inq-four-world-calculations}{Further support and resolution calculations}
Use the four valuations \(00,01,10,11\), with the first coordinate for
\(p\), and put \(P_0=\{00,01\},P_1=\{10,11\}\),
\(Q_0=\{00,10\},Q_1=\{01,11\}\). The following are the maximal supports:
\begin{longtable}{@{}p{0.42\linewidth}p{0.50\linewidth}@{}}
\caption{Four-world support calculations}\label{tab:p08-inq-four-world}\\
\hline Formula & Alternatives \\\hline\endfirsthead
\caption[]{Four-world support calculations (continued)}\\
\hline Formula & Alternatives \\\hline\endhead
\hline\multicolumn{2}{r}{\emph{continued on next page}}\\
\endfoot
\endlastfoot
\(p\land?q\) & \(\{10\},\{11\}\). \\
\(?p\mathbin{\vee\!\vee}?q\) & \(P_0,P_1,Q_0,Q_1\). \\
\(?\neg p\) & \(P_0,P_1\). \\
\(?(p\lor q)\land?(p\land q)\) & \(\{00\},\{01,10\},\{11\}\). \\
\(\neg q\to?p\) & \(\{00,01,11\},\{01,10,11\}\). \\
\((p\lor q)\to(?p\land?q)\) & \(\{00,01\},\{00,10\},\{00,11\}\). \\
\(?p\to?\neg p\) & \(W\). \\
\((?p\to?q)\to?q\) & \(P_0,P_1,Q_0,Q_1\). \\
\hline
\end{longtable}
For syntactic resolutions, \(R(p\land(q\mathbin{\vee\!\vee}r))
=\{p\land q,p\land r\}\), and
\(R(p\to?q)=\{p\to q,p\to\neg q\}\).
The resolutions of \((p\mathbin{\vee\!\vee}q)\to
(r\mathbin{\vee\!\vee}t)\) are
\[
 \{(p\to u)\land(q\to v):u,v\in\{r,t\}\},
\]
where \(p,q,r,t\) are atoms. The resolutions of
\(?(p\lor q)\land?(p\land q)\) are the four conjunctions choosing
one of \(p\lor q,\neg(p\lor q)\) and one of
\(p\land q,\neg(p\land q)\). One is inconsistent, but remains in
the syntactic resolution set. Removing nonmaximal truth sets is an
operation on alternatives, not part of the recursive definition of \(R\).
\end{hbexample}

\begin{hblemma}{thm:p08-inq-finite-calculations}{proved=annex:p8.inq.finite-calculations}{Verification of the finite lookup cases}
The alternative families and entailment tables in
Examples~\ref{def:p08-inq-code-table}, \ref{def:p08-inq-entailment-table}
and \ref{def:p08-inq-four-world-calculations}, and the conditional die
calculations in Example~\ref{def:p08-inq-die-implications}, follow from
the stated support clauses. The syntactic resolution calculations retain
all members, including inconsistent and redundant resolutions.
\end{hblemma}

\section{Semantic scope and neighboring approaches}

The use of all subsets is intentional. Suppose a die outcome can only be
learned completely. If one tested only the initial state and singleton
outcomes, parity would appear to entail range, because every accessible
parity-resolving state already supplies the whole outcome. The state
\(\{2,4\}\) shows that this is not a logical dependency. In a device
revealing only outcome 1, outcome 6, or the block \(\{2,3,4,5\}\), the
same restriction would make even appear to entail high. The singleton
\(\{2\}\) refutes that semantic inference. Reachability in an inquiry
protocol and consequence between information contents answer different questions.

Alternative information semantics can take enhancement states as primitive,
rather than sets of worlds. They make different structural assumptions
explicit and can yield different logics. Similarly, allowing nonpersistent
acceptance conditions for epistemic possibility is a different extension.
Neither variation is silently incorporated into the support clauses above.

Answer-set approaches, partition logic, interrogative inquiry and inferential
erotetic logic retain their own purposes and consequence notions. The generator
construction explains an exact bridge where an answer family generates the
same support proposition; an arbitrary answer list or a question-evocation
relation is not thereby identified with support entailment. The present
proofs use questions as assumptions about information being available,
and conclusions about information being obtainable, without identifying
such proofs with a sequence of actual inquiry actions.

Encoding a question as a modal statement such as \(\Box p\lor\Box\neg p\)
requires stating what information the accessibility relation represents.
If statements are evaluated at a world while questions are evaluated at its
information state, ordinary mixed consequence is not recovered automatically.
Replacing local consequence by truth at every world would, for example,
permit the unwanted move from \(p\) to \(\Box p\). A uniform support
interpretation avoids that change of consequence relation.

The primitive split rule concerns classical antecedents; replacing it by
a Kreisel--Putnam principle is justified in this classical base because
classical formulas are equivalent to negations. That replacement cannot
be transferred without proof to inquisitive logics with nonclassical
statements. First-order and modal inquisitive systems likewise have their
own expressiveness, compactness and axiomatization questions. None of the
propositional completeness or finite-resolution claims here asserts a
corresponding theorem for those larger languages.

\index{inquisitive logic!resolution algorithm}\index{inquisitive logic!canonical support model}
\index{inquisitive logic!expressive completeness}\index{inquisitive logic!coherence degree}
\index{inquisitive logic!legitimate substitution}\index{inquisitive logic!connective independence}
